\chapter{Descrizione delle Strutture Dati}

\section{Edge}

La struttura \texttt{Edge} rappresenta una tratta del grafo nella lista 
di adiacenza. È definita come segue:

\begin{lstlisting}
struct Edge {
    int dest;
    int cost;
};
\end{lstlisting}

\begin{itemize}
    \item \texttt{dest} --- indice della città di destinazione;
    \item \texttt{cost} --- costo intero positivo della tratta.
\end{itemize}

\section{State}

La struttura \texttt{State} rappresenta uno stato nella coda di priorità 
della variante di Dijkstra. A differenza dell'algoritmo classico, che 
mantiene solo la distanza dal nodo sorgente, è necessario tenere traccia 
anche della tratta più costosa percorsa fino a quel momento:

\begin{lstlisting}
struct State {
    int totalCost;
    int maxCost;
    int node;

    int effectiveCost() const {
        return totalCost - maxCost;
    }

    bool operator>(const State& other) const {
        return effectiveCost() > other.effectiveCost();
    }
};
\end{lstlisting}

\begin{itemize}
    \item \texttt{totalCost} --- somma dei costi di tutte le tratte 
    percorse dal nodo sorgente fino al nodo corrente;
    \item \texttt{maxCost} --- costo della tratta più onerosa nel 
    percorso corrente, che rappresenta lo sconto applicabile;
    \item \texttt{node} --- indice del nodo corrente;
    \item \texttt{effectiveCost()} --- calcola il costo effettivo dopo 
    lo sconto: $\texttt{totalCost} - \texttt{maxCost}$;
    \item \texttt{operator>} --- operatore di confronto per la 
    \texttt{priority\_queue}, che ordina gli stati per costo effettivo 
    minimo (min-heap).
\end{itemize}

\section{Graph}

La classe \texttt{Graph} incapsula l'intera struttura del grafo e 
le operazioni su di esso. Internamente mantiene:

\begin{itemize}
    \item \texttt{int N} --- numero di città nel grafo;
    \item \texttt{vector<vector<Edge>> adj} --- lista di adiacenza: 
    per ogni città $v$, \texttt{adj[v]} contiene la lista delle tratte 
    uscenti da $v$.
\end{itemize}

\subsection{Lista di Adiacenza}

La rappresentazione scelta è la \textbf{lista di adiacenza} tramite 
\texttt{std::vector<std::vector<Edge>>}. Questa scelta è motivata dalle 
caratteristiche del problema:

\begin{itemize}
    \item il numero di tratte $P$ può essere molto inferiore a $N^2$, 
    rendendo la matrice di adiacenza spreca di memoria;
    \item la lista di adiacenza occupa $O(N + P)$ in memoria;
    \item l'iterazione sulle tratte uscenti da un nodo, operazione 
    centrale in Dijkstra, è efficiente con questa rappresentazione.
\end{itemize}

\subsection{Vettori Ausiliari in Dijkstra}

Durante l'esecuzione di \texttt{dijkstra()} vengono utilizzati i 
seguenti vettori ausiliari di dimensione $N$:

\begin{itemize}
    \item \texttt{dist[v]} --- miglior costo effettivo trovato per 
    raggiungere la città $v$;
    \item \texttt{totalCost[v]} --- costo totale del percorso ottimo 
    verso $v$, necessario per l'output;
    \item \texttt{maxEdge[v]} --- tratta più costosa nel percorso ottimo 
    verso $v$, necessaria per l'output;
    \item \texttt{parent[v]} --- città precedente nel percorso ottimo 
    verso $v$, necessario per la ricostruzione del percorso.
\end{itemize}